\section{Finite-time approximation under a global regulator}\label{sec:finite}
The two endpoint formulas do not yet provide informative overlap through
the intervening times. We now fix one smooth measurement protocol and
construct a finite heat field on any compact time slab. Its error is
controlled at the source, then transferred to the density, current and
actual covariance. The construction has a finite explicit cost; the
bounds alone do not assert that this cost is small.

\subsection{One regulator for both endpoints}
Let $C=\Phi(0)\sqrt\pi/2$, fix $\eta>0$, and define
\[
 \epsilon_b(t)=\eta C(1+t^2)^{-1/4}.
\]
For $64<t<128$ put $v=(t-64)/64$ and
$\chi(t)=e^{-1/v}/(e^{-1/v}+e^{-1/(1-v)})$; extend $\chi$ by zero
for $t\le64$ and one for $t\ge128$. Set
\begin{equation}
 \epsilon_g(t,x)=
 \begin{cases}
 \epsilon_b(t),&t\le64,\\
 \exp\{(1-\chi(t))\log\epsilon_b(t)+\chi(t)\log\epsilon_+(t,x)\},
       &t>64.
 \end{cases}
 \label{eq:globalRegulator}
\end{equation}
This is positive and smooth at every finite $(t,x)$.
To check that the large saddle used in $\epsilon_+$ exists throughout
the switch, the maximum of $q_t'$ is
$(t/2)(\log(t/(8\pi))-1)+9$, attained at
$u_*=\log(t/(8\pi))/4$.
It is positive at $t=64$ and increasing thereafter; the larger zero
has $\kappa>0$. Smooth extension slightly below $64$ and the flatness
of $\chi$ justify the gluing. The switch points specify the regulator;
they are not onset times for Theorem~\ref{thm:explicit}.

At $t=-T$,
$\epsilon_b/ (\eta C_T)=(1+T^{-2})^{-1/4}=1+O(T^{-2})$.
Using \eqref{eq:Eminus} with this regulator and replacing $\eta$ by
$\eta(1+T^{-2})^{-1/4}$ in \eqref{eq:adaptiveNegative} preserves its
$O(T^{-1})$ scaled error. Indeed
$|\partial_{\log\eta}F_\eta|\le2F_\eta$, and $F_\eta$ is bounded on
each fixed scaled transverse strip, uniformly over real scaled heights.
For $t\ge128$, \eqref{eq:globalRegulator} equals
\eqref{eq:positiveRegulator} exactly.
On a compact time and spatial buffer it has a positive lower bound.
After the height dependence turns on, positivity alone does not supply
such a bound over all real heights.

\subsection{A finite source field with error-function primitives}
Choose a time slab $[a,b]$, transverse buffer $|y|\le Y$, cutoff $U>0$,
number of theta modes $N\ge1$, and nodes
$0=u_0<\cdots<u_m=U$ with largest step $h_s$.
Let $\Phi_N$ be the first $N$ terms of \eqref{eq:phi} and let $f$
be its continuous piecewise linear interpolant on $[0,U]$, extended
by zero for $u>U$. Each nodal value is positive, so $f\ge0$.
The possible jump at $U$ is allowed. Define
\begin{equation}
 H_f(t,z)=\int_0^U e^{tu^2}f(u)\cos(zu)\,du.
 \label{eq:finitefield}
\end{equation}
This integral has a finite error-function expression.
On each interval write $f(u)=A_i+B_i u$ and use
\begin{align}
 I_0(t,z;u)&=\frac{\sqrt\pi}{2\sqrt{-t}}e^{z^2/(4t)}
 \left[\operatorname{erf}\left(\sqrt{-t}\,u-\frac{iz}{2\sqrt{-t}}\right)
       -\operatorname{erf}\left(-\frac{iz}{2\sqrt{-t}}\right)\right],
 \label{eq:primitive0}\\
 I_1(t,z;u)&=\frac{e^{tu^2+izu}-1-izI_0(t,z;u)}{2t},\qquad
 C_j(t,z;u)=\frac{I_j(t,z;u)+I_j(t,-z;u)}2.
 \label{eq:primitive1}
\end{align}
Then $H_f=\sum_i(A_i\Delta_iC_0+B_i\Delta_iC_1)$, where
$\Delta_iC_j=C_j(u_{i+1})-C_j(u_i)$.
The branch choice in the square root does not change the result.
At $t=0$ the removable limits are
\[
 I_0=(e^{izu}-1)/(iz),\qquad
 I_1=\{e^{izu}(izu-1)+1\}/(iz)^2,
\]
with values $u$ and $u^2/2$ when also $z=0$.
Completing the square and differentiating in $u$ prove the formulas.
The field is entire in $(t,z)$, even in $z$, real on the real axis,
and satisfies $(H_f)_t=-(H_f)_{zz}$ exactly.
Its coefficients come from prescribed source values, without fitting
theta roots or densities. Numerical cancellation in the primitives
requires its own evaluation allowance.

\subsection{A complete source-to-jet error}
Let $B_2$ be the source derivative bound of Lemma~\ref{lem:B2} and put
\begin{equation}
 \tau_N=\frac{2\pi^2(N+1)^4e^{-\pi(N+1)^2}}{1-\vartheta_N},\qquad
 \vartheta_N=\left(1+\frac1{N+1}\right)^4e^{-\pi(2N+3)},\qquad
 d_s=B_2h_s^2/8+\tau_N.
 \label{eq:sourceinterpolation}
\end{equation}
For $s=\max\{b,0\}$ choose $U$ so that, for $0\le j\le3$,
\begin{equation}
 16\pi e^{4U}\ge2s,\qquad
 \lambda_j=4\pi e^{4U}-2sU-9-Y-j/U>0.
 \label{eq:tailgates}
\end{equation}
Define
\begin{align}
 T_j&=\frac{2\pi^2C_bU^j
               e^{sU^2+(9+Y)U-\pi e^{4U}}}{\lambda_j},\notag\\
 F_j&=\int_0^U u^je^{su^2+Yu}\,du
       \le\frac{U^{j+1}}{j+1}e^{sU^2+YU},\qquad
 \delta_j=d_sF_j+T_j.\label{eq:jetallowances}
\end{align}
The $F_j$ also have finite error-function recurrences.

\begin{theorem}[Uniform finite-time jet approximation]\label{thm:finitejets}
Under \eqref{eq:tailgates}, for all $t\in[a,b]$, all real $x$ and
$|y|\le Y$,
\begin{equation}
 |\partial_z^j(H-H_f)(t,x+iy)|\le\delta_j,\qquad 0\le j\le3.
 \label{eq:finitejets}
\end{equation}
For every fixed finite slab and strip, arbitrary positive jet tolerances
are attainable by first enlarging $U$, then $N$ and the mesh resolution.
\end{theorem}
\begin{proof}
Each omitted positive leading theta term decreases on $u\ge0$;
its value at zero and the consecutive mode ratio give
$0\le\Phi-\Phi_N\le\tau_N$.
Interpolation of the complete $\Phi$ costs at most $B_2h_s^2/8$,
and interpolation of its finite nodal approximations adds at most
$\tau_N$. Thus $|\Phi-f|\le d_s$ on $[0,U]$.
On $[U,\infty)$ the logarithm of
$u^je^{su^2+(9+Y)u-\pi e^{4u}}$ is concave by
\eqref{eq:tailgates} and has derivative $-\lambda_j$ at $U$.
Its tangent exponential integrates to the tail bound $T_j$.
Finally $|\partial_z^j\cos(zu)|\le u^je^{Y u}$, independently of $x$,
so the two source regions contribute $d_sF_j$ and $T_j$.
For fixed $s,Y$, the tails tend to zero with $U$; after fixing $U$,
the finite-mode and interpolation errors tend to zero with $N$ and
mesh refinement.
\end{proof}
The heat equation also pays the corresponding time derivatives:
$\partial_t\partial_z^j(H-H_f)=-\partial_z^{j+2}(H-H_f)$.
Nodal arithmetic error $d_{\rm num}$ adds $d_{\rm num}F_j$ to
$\delta_j$. Primitive evaluation and spatial averaging errors remain
separate. Refining the mesh alone cannot remove a fixed mode or tail
allowance.

\subsection{Density and transport stability through zeros}
Fix a full comparison buffer where $\epsilon_g\ge e_*>0$ and
$|r_t|\le a_0$. Both the true and approximate jets are bounded by
\begin{equation}
 P_j=\int_0^U u^je^{su^2+Yu}f(u)\,du+\delta_j.
 \label{eq:finitejetmajorants}
\end{equation}
These positive integrals are finite combinations of the same primitives.
Let $c_0=125/(54\sqrt5)$ and define
\begin{align}
 E_\rho&=\frac{2P_1\delta_1}{\pi e_*^2}
                   +\frac{c_0P_1^2\delta_0}{\pi e_*^3},\notag\\
 E_j&=\frac{P_1\delta_2+P_2\delta_1}{\pi e_*^2}
                   +\frac{c_0P_1P_2\delta_0}{\pi e_*^3},\notag\\
 E_3&=\frac{P_1\delta_3+P_3\delta_1}{\pi e_*^2}
                   +\frac{c_0P_1P_3\delta_0}{\pi e_*^3},\qquad
 \rho_{\max}=P_1^2/(\pi e_*^2),\label{eq:fieldstability}\\
 E_{\rho,t}&=2E_3+a_0(E_\rho+2\rho_{\max}\delta_0/e_*)
       +2P_2E_\rho/e_*+2\rho_{\max}\delta_2/e_*
       +4\rho_{\max}P_2\delta_0/e_*^2.\notag
\end{align}

\begin{proposition}[Regulated field stability]\label{prop:fieldstability}
Use the same regulator at each point for $H$ and $H_f$. Then
\begin{equation}
 |\rho[H]-\rho[H_f]|\le E_\rho,\quad
 |j[H]-j[H_f]|\le E_j,\quad
 |\rho_t[H]-\rho_t[H_f]|\le E_{\rho,t}.
 \label{eq:stabilityconclusions}
\end{equation}
Positive normalized convolution preserves these sup-norm bounds.
No simple-zero premise is needed, so the comparison includes multiple
zeros and nonzero stationary points.
\end{proposition}
\begin{proof}
For $A(h)=q/[\pi(|h|^2+q)^2]$,
$A\le1/(\pi e_*^2)$ and the global gradient bound is
$|\nabla_h A|\le c_0/(\pi e_*^3)$, as in
Theorem~\ref{thm:global}. Product comparison proves the first two
bounds and the analogous product bound $E_3$.
Writing $p=H_z$, $v=H_{zz}$ and $w=q/(|H|^2+q)$, the exact derivative is
\[
 \rho_t=-2A\Re(H_{zzz}\overline p)+r_t(1-2w)\rho
              +4\rho\Re\frac{v\overline H}{|H|^2+q}.
\]
Use $|w[H]-w[H_f]|\le\delta_0/e_*$,
$|H|/(|H|^2+q)\le1/(2e_*)$, and the Lipschitz bound
$\delta_0/e_*^2$ for $H/(|H|^2+q)$. Termwise comparison gives exactly
$E_{\rho,t}$. This compares the actual time derivatives and does not
differentiate the static error allowance.
\end{proof}
For a slab ending at $b\le64$, $\epsilon_g$ depends only on time,
so these estimates are uniform over all real heights. On later slabs
a finite horizontal buffer supplies $e_*$ and $a_0$.

\subsection{The actual covariance and its Laplacian}
Let $S_f=\log(|H_f|^2+q)$, $U_f=K*S_f$ and
\[
 Q_f=K*((S_f)_x^2-(S_f)_y^2)-((U_f)_x^2-(U_f)_y^2).
\]
This is computed from the approximate field, not chosen to force a
potential equation. For $i=x,y$ put
\begin{equation}
 b_i=\|r_i\|_\infty,\qquad
 G_i=P_1/e_*+b_i,\qquad
 d_i=\delta_1/e_*+2P_1\delta_0/e_*^2+b_i\delta_0/e_*.
 \label{eq:gradientstability}
\end{equation}
Then
\[
 |S-S_f|\le\delta_0/e_*,\qquad
 |S_i|,|(S_f)_i|\le G_i,\qquad |S_i-(S_f)_i|\le d_i.
\]
These follow from
$S_i=2\Re(H_i\overline H/(|H|^2+q))+wr_i$ and the same bounded
rational maps used above.

\begin{theorem}[Filtered covariance stability]\label{thm:covariance}
Let $K$ be fixed, nonnegative, normalized and compactly supported,
with the full comparison buffer supplied. Set
$L_1=\|K_x\|_1+\|K_y\|_1$ and $L_2=\|\Delta K\|_1$.
Then
\begin{equation}
 \|Q-Q_f\|_\infty\le4\sum_{i=x,y}G_i d_i,\qquad
 \|\Delta(Q-Q_f)\|_\infty
       \le(4L_1^2+6L_2)\sum_{i=x,y}G_i d_i.
 \label{eq:covariancestability}
\end{equation}
\end{theorem}
\begin{proof}
The difference of the averaged gradient squares costs
$2\sum_iG_i d_i$, as does the difference of the squared means.
For the Laplacian, move both derivatives of the first averaged squares
onto $K$, costing $2L_2\sum_iG_i d_i$.
For a mean gradient $v_i=K*S_i$ use
$\Delta(v_i^2)=2v_i\Delta v_i+2|\nabla v_i|^2$.
Comparing true and approximate mean gradients costs
$4(L_2+L_1^2)\sum_iG_i d_i$. Adding the two costs proves the result.
The $C^3$ kernel \eqref{eq:kernel} has all the required integrable
derivatives. No pointwise third spatial gradient of $S$ is required.
\end{proof}
The potential bound also controls its convolved boundary traces wherever
the full buffer is provided; an arbitrary harmonic correction has not
been introduced. Since $H_f$ solves heat exactly, its density, current,
potential and covariance satisfy Proposition~\ref{prop:adaptive} with
all three regulator corrections retained.
These are source-based comparisons with the true quantities. They do
not give an autonomous covariance closure in $R_f$ alone or transfer
the clock's transport error to an unknown clock covariance.

\subsection{Source locality removes artificial cutoff amplification}
The global bound \eqref{eq:jetallowances} loses the source decay before
applying the heat weight. A cellwise replacement preserves it. Let
$[a_i,b_i]$ be a source cell of length $h_i$, let $B_i$ bound
$\sup_{[a_i,b_i]}|\Phi''|$, and let $\lambda_i$ be the nodal hats.
For implemented nodal data $v_i$, assume
$|v_i-\Phi_N(u_i)|\le\sigma_i$, with nonnegative interpolant $f$.
Define the spatially dependent omitted-mode allowance
\begin{equation}
 \tau_N(v)=\frac{2\pi^2(N+1)^4e^{9v-\pi(N+1)^2e^{4v}}}
 {1-(1+1/(N+1))^4e^{-\pi(2N+3)e^{4v}}}.
 \label{eq:localmodetail}
\end{equation}
For a time ceiling $b$, write $W_j(u)=u^je^{bu^2+Yu}$ and retain
$T_j$ from \eqref{eq:jetallowances} with $s=\max\{b,0\}$.

\begin{proposition}[Cellwise jet comparison]\label{prop:localjets}
Uniformly for $t\le b$ in the chosen slab, all real $x$, and $|y|\le Y$,
\begin{align}
 |\partial_z^j(H-H_f)|&\le\delta_j^{\rm loc},\qquad 0\le j\le3,\notag\\
 \delta_j^{\rm loc}
 &=\frac12\sum_i B_i\int_{a_i}^{b_i}(u-a_i)(b_i-u)W_j(u)\,du\notag\\
 &\quad+\sum_i\bigl(\tau_N(u_i)+\sigma_i\bigr)
                \int_0^U\lambda_i(u)W_j(u)\,du+T_j.
 \label{eq:localjets}
\end{align}
For the complete-source cell constants below, the interpolation term is
no larger than $B_2^{\Phi}h_s^2F_j/8$. For fixed $b,Y,j$ and maximum cell length
bounded above, its coefficient of $h_s^2$ remains bounded as $U$ grows.
\end{proposition}
\begin{proof}
Linear interpolation of the complete source costs
$B_i(u-a_i)(b_i-u)/2$. The positive hats transfer the nodal mode and
arithmetic errors, while consecutive omitted-mode ratios prove
\eqref{eq:localmodetail}. Integrating these bounds against $W_j$ and
adding the complete source tail proves \eqref{eq:localjets}.

To specify $B_i$, put $A_n=\pi n^2e^{4a_i}$ and $B_n=\pi n^2e^{4b_i}$.
Assign coefficients $75,330,224,32$ to the powers $5/4,9/4,13/4,17/4$,
respectively. Write $m_p(A,B)=v^pe^{-v}$, where
$v=\min\{B,\max\{A,p\}\}$. Then
\[
 B_i=\pi^{-1/4}\sum_{n\ge1}n^{-1/2}\sum_p c_p m_p(A_n,B_n)
 \le B_2^{\Phi}.
\]
This follows by bounding the four monomials in the differentiated
source as in Lemma~\ref{lem:B2}. For a finite sum through $N_B\ge1$,
the omitted tail is at most its first omitted expression at $a_i$
divided by
$1-(1+1/(N_B+1))^8e^{-\pi(2N_B+3)e^{4a_i}}$.
Each omitted monomial is decreasing in $u$.

For cell lengths at most $h_0$, a bound for $B_i$ at $u$ is the
supremum of this positive derivative envelope over
$|v-u|\le h_0$, $v\ge0$. Its far tail is bounded by a polynomial in
$e^{4(u+h_0)}$ times $e^{-\pi e^{4(u-h_0)}}$. This is integrable
against $W_j$. Using $(u-a_i)(b_i-u)\le h_s^2/4$ proves the final
assertion. The original coefficient $B_2^{\Phi}F_j/8$, in contrast, diverges
with $U$ when $b>0$. This distinction concerns artificial cutoff
amplification, not the genuine time growth of the approximation cost.
\end{proof}
A quadrature-free upper bound for each interpolation contribution is
$B_i h_i^3\max_{[a_i,b_i]}W_j/12$; the nodal contribution is at most
$h_i\max W_j$ times the mean of its two endpoint allowances.
The positive maxima occur at endpoints or positive solutions of
$2bu^2+Yu+j=0$. Taking the minimum with the original bound preserves
dominance even on a coarse mesh. Nodal arithmetic retains the source
decay only when its pointwise allowances do; a relative rounding
hypothesis must be verified. Note~23 \cite{n23} records the separate
interpolation, mode, tail and arithmetic terms.

\subsection{Local amplitude improves nonlinear transfer}
At a comparison point put $g=H_f$, $g_j=\partial_z^jH_f$, and let
$d_j$ be any valid jet allowances. Set
\[
 P_j=|g_j|+d_j,\quad m=(|g|-d_0)_+,\quad M=|g|+d_0,\quad q=\epsilon_g^2.
\]
Define maxima over $m\le v\le M$:
\begin{align}
 A_*&=\frac{q}{\pi(q+m^2)^2},&
 A_1&=\max_v\frac{4qv}{\pi(q+v^2)^3},\notag\\
 a_*&=\max_v\frac{v}{q+v^2},& b_*&=\frac1{q+m^2},&
 w_1&=\max_v\frac{2qv}{(q+v^2)^2}.\label{eq:localrational}
\end{align}
The three maxima are evaluated at the clamped points
$\epsilon_g/\sqrt5$, $\epsilon_g$, and $\epsilon_g/\sqrt3$.
Set
\begin{align}
 e_\rho&=2A_*P_1d_1+A_1P_1^2d_0,\notag\\
 e_j&=A_*(P_1d_2+P_2d_1)+A_1P_1P_2d_0,\notag\\
 e_{13}&=A_*(P_1d_3+P_3d_1)+A_1P_1P_3d_0,\notag\\
 e_w&=w_1d_0,\qquad e_S=2a_*d_0,\qquad \rho_*=A_*P_1^2,\notag\\
 e_{\rho,t}&=2e_{13}+|r_t|(e_\rho+2\rho_*e_w)
 +4\{P_2a_*e_\rho+\rho_*(a_*d_2+P_2b_*d_0)\}.
 \label{eq:localnonlinear}
\end{align}

\begin{proposition}[Amplitude-aware comparison]\label{prop:localnonlinear}
The quantities in \eqref{eq:localnonlinear} bound the corresponding
actual density, current, potential and density-time-derivative errors.
In addition,
\begin{align}
 |S_t-(S_f)_t|&\le2(a_*d_2+P_2b_*d_0)+|r_t|e_w,\notag\\
 |S_i|,|(S_f)_i|&\le G_i:=2P_1a_*+|r_i|,\notag\\
 |S_i-(S_f)_i|&\le d_i^S:=2a_*d_1+2P_1b_*d_0+|r_i|e_w.
 \label{eq:localpotential}
\end{align}
No simple-root or nonvanishing-field premise is needed.
\end{proposition}
\begin{proof}
The segment between the two field values lies in the annulus
$m\le|h|\le M$. The constants \eqref{eq:localrational} bound
$A(h)$, $|\nabla_h A|$, $|h/(|h|^2+q)|$, the derivative norm of
that map, and $|\nabla_h w|$, respectively.
For the fourth bound the radial eigenvalues are
$(q-v^2)/(q+v^2)^2$ and $1/(q+v^2)$.
The product comparison and exact time-derivative argument from
Proposition~\ref{prop:fieldstability} now give
\eqref{eq:localnonlinear}. The exact formulas for $S_t,S_i$ give
\eqref{eq:localpotential} directly.
All constants are bounded by their global counterparts in
\eqref{eq:fieldstability}--\eqref{eq:gradientstability}.
For example, relative to $1/(\pi\epsilon_g^2)$, $A_*$ gains the
factor $[1+(m/\epsilon_g)^2]^{-2}$.
The gain may disappear at a zero; the bound remains valid there.
\end{proof}
No physical normalization has changed. In particular there is no
unpaid time derivative of a rescaled field.

\subsection{Local covariance and evaluation budgets}
The local envelopes may be filtered before taking suprema.
For any pointwise error $e_V$,
$|\partial^\alpha K*(V-V_f)|\le|\partial^\alpha K|*e_V$.
With $d_i=d_i^S$ from \eqref{eq:localpotential}, direct comparison of
the actual covariances gives
\begin{align}
 |Q-Q_f|&\le2\sum_i\{K*(G_id_i)+(K*G_i)(K*d_i)\},\notag\\
 |\Delta(Q-Q_f)|&\le\sum_i\biggl\{2|\Delta K|*(G_id_i)
 +2(K*d_i)(|\Delta K|*G_i)\notag\\
 &\qquad+2(K*G_i)(|\Delta K|*d_i)
 +4\sum_{a=x,y}(|K_a|*G_i)(|K_a|*d_i)\biggr\}.
 \label{eq:localcovariance}
\end{align}
To prove this, move the derivatives of the averaged squares onto $K$
and use $\Delta(v^2)=2v\Delta v+2|\nabla v|^2$ for each averaged
gradient. This is the proof of Theorem~\ref{thm:covariance} with its
spatial information retained. For constant envelopes the Laplacian
coefficient becomes
$6\|\Delta K\|_1+4(\|K_x\|_1^2+\|K_y\|_1^2)$, slightly sharper
than \eqref{eq:covariancestability}. Positive convolution bounds still
need a numerical integration enclosure if evaluated by quadrature.

There is also a finite evaluation formula without division by $t$ or $z$.
On a cell centered at $c$ with half-width $d$, expand
\[
 e^{tu^2\pm izu}=e^{tc^2\pm izc}
 \sum_{n=0}^{M_e}\frac{[(2tc\pm iz)(u-c)+t(u-c)^2]^n}{n!}
       +\text{remainder}.
\]
Multiply by the linear source and $u^j$, and integrate the polynomial
exactly, with derivative factors $(\pm i)^j/2$ for the two branches.
Writing $b=c+d$ and $\Theta=(2|t|c+|z|)d+|t|d^2$, the cell's
absolute jet truncation error is at most
\begin{equation}
 2d\max\{f(c-d),f(c+d)\}b^j
 e^{tc^2+|y|c+\Theta}\frac{\Theta^{M_e+1}}{(M_e+1)!}.
 \label{eq:localevaluation}
\end{equation}
This follows from the absolute exponential remainder. Sum over cells.
It pays truncation, not rounding. The exact $H_f$ still solves heat;
a truncated implementation is compared to it with this additional error.
The raw error-function subtraction can be unusable near its removable
singularity even at high precision; Note~23 records a failed case and
the separate local evaluation check.

For the first trial, Note~23 fixes $[-16,64]$, observation rectangle
$[-8,8]\times[-1/2,1/2]$, the full buffer
$[-10,10]\times[-1,1]$, widths $(2,1/2)$ and $\eta=1/2$.
Its absolute contract includes $R,J,R_t,U,Q,\Delta Q$, the required
filtered derivatives and source-derived boundary traces.
On this slab $r_x=r_y=0$, so the spatial curvature correction is exactly
zero, while the time forcing and current source remain.
The 192-cell diagnostic does not meet the contract. The proved local
bounds are an improvement to the comparison method, not an informative
full-transition representation or a certified numerical solution.


\subsection{What the finite bridge does and does not complete}
The construction proves approximability on each compact slab.
A useful transition bridge must still choose and evaluate $U,N$, the
mesh, widths, $\eta$ and the time interval, and show informative overlap
with both endpoint formulas using the complete error budgets.
Where two independently paid approximations overlap, the triangle
inequality bounds their difference. A time blend of heat fields has
an additional heat residual: $(1-\chi)H_f+\chi G$ contributes
$\chi'(G-H_f)$ even if each component solves heat. No such blend is
asserted here.

\begin{proposition}[Failure of a fixed compact source at positive infinity]
\label{prop:fixedsource}
Fix $U$ and $f$ in \eqref{eq:finitefield}. Under the regulator
\eqref{eq:positiveRegulator}, its density tends to zero uniformly on
fixed spatial buffers as $t\to\infty$. For the true theta field, at
any $y$ with $k_h(y)>0$, the corresponding averaged density grows
like $Lk_h(y)/\pi$. Thus this fixed compact-source field cannot supply
the positive endpoint of an all-time approximation.
\end{proposition}
\begin{proof}
For fixed $U,Y$, each fixed jet of $H_f$ is $O(e^{tU^2+YU})$.
But $\log M_0\sim tL^2$ and $L\to\infty$, so $\epsilon_+$
outgrows every $e^{tU^2}$ on a fixed spatial buffer.
The bound $\rho[H_f]\le|(H_f)_z|^2/(\pi\epsilon_+^2)$ proves
vanishing. For theta, Theorem~\ref{thm:explicit} and the periodization
bound give $K*\rho=Lk_h/\pi+o(1)$ there.
\end{proof}
This exclusion is specific to a fixed compact source; it does not
exclude growing cutoffs. It also shows why an exactly zero heat
residual alone cannot establish approximation to theta.

The remaining analytical tasks are informative transition overlap,
a positive outer-height model and a potential/covariance comparison for
the explicit clock on a domain suitable for evolution.
Certified finite onset and independent mathematical review remain
separate tasks. The endpoint asymptotics, determinant comparisons and
finite-slab stability inequalities do not exclude earlier collisions
or establish simplicity at a prescribed earlier time.
